% 04_setup -- written by Claude Code (D-045); revised after the peer review (Prompt 16: M1, M6).
\section{Experimental Setup}\label{sec:setup}
Generation ran on Kaggle's free tier with two NVIDIA T4 GPUs (16\,GB each), using vLLM 0.32.0~\cite{kwon2023vllm}
with PyTorch 2.13, tensor parallelism over both GPUs and a context length of 2{,}048 tokens. Both arms were generated
in batches of 64 items, so one engine call held 64 sequences for greedy decoding and 640 for sampling; GPU time
therefore depends on this configuration, and we use output tokens as the measure of cost. Generation took
\genHoursMain{}~h for the main run and \genHoursExtras{}~h for the exploratory run (\sessionHoursMain{} and
\sessionHoursExtras{}~h of GPU session). Each notebook installed pinned versions, cloned the repository at a fixed
commit, checked dataset checksums and ran the unit tests before generating. Evaluation ran on a CPU (Python 3.12,
pandas 3.0, NumPy 1.26, SciPy 1.17), after an audit of the saved outputs (split checksums, no question shared between
debug and test sets, prompts independent of the answer key, identical items, engine, model revision, seeds and token
limit in every arm, chance accuracy with shuffled keys; \auditTotal{} of \auditTotal{} checks passed). Re-running the
evaluation on another machine reproduced the Kaggle result files byte for byte. Greedy decoding on GPUs is not exactly
deterministic (\rerunFiftySame{} of 50 letters were the same in the 50-item rerun) and sampled outputs are not
bit-reproducible. The raw outputs restate licensed question text and are not released; every statistic can be
recomputed exactly from the released parsed answers, and every number in the text is inserted by a script.
